\clearpage

\begin{center}
{\Large \textbf{Abstract}}
\end{center}

\vspace{1cm}

This thesis studies problems related to matchings in graphs. Specifically, in the first chapters, we study two classical problems, namely finding a maximum matching in bipartite graphs and the stable marriage problem. We approach these two problems first from an algorithmic perspective and then from the perspective of linear and integer programming. In addition, we study the house allocation problem under the criterion of Pareto optimality. We show that the linear relaxation of the corresponding integer program is not integral and solve the separation problem in $O(n^{3}L^{3})$ time, where $n$ is the number of applicants and $L$ is the maximum length of a preference list. In particular, when there are $n$ applicants and $n$ houses, this is equal to $O(n^{6})$. We also conduct computational experiments on small instances by constructing the convex hull of the vectors corresponding to Pareto optimal matchings and comparing it with the linear relaxation. Furthermore, we prove the validity of two families of additional constraints and show that their inclusion strengthens the linear relaxation. For a cyclic parametric family of house allocation instances, we also derive a new family of cycle-break inequalities and prove that they are facet-defining for the corresponding Pareto optimal matching polytope. We also prove that the dimension of the polytope of Pareto optimal matchings is at most $|E|-|F|$, where $|F|$ is the number of houses that appear as the first choice of at least one applicant. Finally, we present and describe the source code developed for conducting the experiments, which contributed to obtaining the above results.

\vspace{2em}
\noindent

\textbf{Keywords:} Matching Theory, Pareto Optimality, Graph Algorithms, Combinatorial Optimization, Polyhedral Methods, Linear Programming, Integer Programming.

\clearpage